% !PRINTED 239
\weberpage{239}
For this it is required that
\weq{11}{\begin{gathered}
\alpha a_2 + \beta c_2 = 0 \\
\gamma a_2 + \delta c_2 = -\partial,
\end{gathered}}
\weq{12}{\begin{gathered}
\beta \partial_2 = a \\
\delta \partial_2 = c.
\end{gathered}}
Thus $\partial_2$ is determined as the greatest common
divisor of $a$ and $c$, and with it at the same time, because of $a_2\partial_2 = n$, also $a_2$.

By the two equations (12) one now knows $\beta$, $\delta$
as relatively prime numbers, and can determine $\alpha$, $\gamma$ from the Diophantine
equation
\[
\alpha\delta - \beta\gamma = 1
\]
% ["bestimmen." closes the German sentence after the display; English requires the verb before it]
When this has been done, there follows from the two
equations (11)
\weq{13}{\begin{gathered}
a_2 = \partial\beta, \\
c_2 = -\partial\alpha,
\end{gathered}}
whereby $c_2$ also is determined, and it appears at the same time that $\partial$
is the greatest common divisor of $a_2$, $c_2$. If one replaces
$\alpha$, $\gamma$ by another solution $\alpha + h\beta$, $\gamma + h\delta$, then only
$c_2$ is altered, by a multiple of $a_2$.

Through the transformation (7), therefore,
\[
j\left(\frac{c+\partial\omega}{a}\right) \text{ is carried to } j\left(\frac{c_2+\partial_2\omega}{a_2}\right)
\]
% ["über." — the separable prefix that closes the German sentence is carried into the display's \text{}]

Let us now form a \gesperrt{symmetric} function of all the
quantities (1), say, for an indeterminate $x$, the product
\[
\Pi\left[x - j\left(\frac{c+\partial\omega}{a}\right)\right],
\]
then this function is not altered by the transformations
(6), (7), and is therefore, by the theorem mentioned above, a \gesperrt{rational
function} of $j(\omega)$. Moreover it is a polynomial
in $x$ of degree $\nu$ with leading term $x^{\nu}$, and
we denote it by $F_n[x, j(\omega)]$. The equation
\weq{14}{F_n[x, j(\omega)] = 0}
has the quantities $j\left(\dfrac{c+\partial\omega}{a}\right)$ as roots, and is the \gesperrt{invariant equation},
whose most important properties we now propose to derive.

1. The invariant equation is \gesperrt{irreducible} when, as domain of rationality, the totality of all rational functions of
%CONT>>
